Training is only half of the energy bill. We deploy each trained model with the confidence-gated cascade of Section~\ref{sec:lifecycle-model}: a single threshold $\theta$ per model is chosen on validation data as the cheapest one whose validation accuracy is within one percentage point of the model's best single exit, and we then report test accuracy and modelled inference energy (fleet-average device mix). Table~\ref{tab:life} compares three deployments: single-exit FedAvg (trained with a last-exit loss; every window runs the full network), multi-exit FedAvg with the cascade, and EcoFed with the cascade. For a fleet that classifies one 2.56~s window every 1.28~s per device, always on, the last column gives the time after which the energy spent on serving exceeds the energy spent on training ($M_\times$ divided by the fleet's window rate).

\begin{table}[t]
\centering\scriptsize
\caption{Lifecycle energy (means over 16 seeds). Test acc.: accuracy of the deployed cascade (single exit for the first row). $E_{train}$: fleet energy used for training. $\bar e_{inf}$: expected inference energy per window. $M_\times=E_{train}/\bar e_{inf}$: number of inferences after which serving costs as much as training did.}\label{tab:life}
\setlength{\tabcolsep}{3pt}
\resizebox{\textwidth}{!}{\begin{tabular}{@{}llrrrrrr@{}}
\toprule
Data / $\phi$ & Model & Test acc. (\%) & $E_{train}$ (J) & $\bar e_{inf}$ ($\mu$J) & exits at 1 (\%) & $M_\times$ ($10^5$) & days to $M_\times$\\
\midrule
UCI HAR, 0.25 & FedAvg, single exit & 87.2 & 116.1 & 845 & 0 & 1.37 & 0.13\\
 & FedAvg, cascade & 83.9 & 116.1 & 250 & 31 & 4.65 & 0.43\\
 & EcoFed, cascade & 85.6 & 70.7 & 114 & 98 & 6.19 & 0.57\\
\midrule
UCI HAR, 1.0 & FedAvg, single exit & 89.1 & 378.9 & 845 & 0 & 4.49 & 0.42\\
 & FedAvg, cascade & 86.5 & 378.9 & 209 & 48 & 18.16 & 1.68\\
 & EcoFed, cascade & 86.4 & 176.5 & 169 & 86 & 10.47 & 0.97\\
\midrule
PAMAP2, 0.25 & FedAvg, single exit & 31.4 & 136.2 & 950 & 0 & 1.43 & 0.11\\
 & FedAvg, cascade & 32.7 & 136.2 & 611 & 7 & 2.23 & 0.17\\
 & EcoFed, cascade & 54.5 & 148.2 & 378 & 73 & 3.92 & 0.29\\
\midrule
PAMAP2, 1.0 & FedAvg, single exit & 47.6 & 555.9 & 950 & 0 & 5.85 & 0.43\\
 & FedAvg, cascade & 46.6 & 555.9 & 450 & 22 & 12.36 & 0.92\\
 & EcoFed, cascade & 58.6 & 472.8 & 436 & 55 & 10.85 & 0.80\\
\bottomrule
\end{tabular}}
\end{table}

Three observations follow. First, \emph{inference energy per window can be cut by a factor of 2 to 7 by the cascade.} EcoFed's cascade needs 114 and 169~$\mu$J per window on UCI HAR against 845~$\mu$J for the single-exit network (7.4 and 5.0 times less), at a cost of 1.6 and 2.7 pp of accuracy at $\phi=0.25$ and $1$; on PAMAP2 it needs 378 and 436~$\mu$J against 950~$\mu$J (2.5 and 2.2 times less) with a \emph{higher} accuracy (54.5\% and 58.6\% against 31.4\% and 47.6\% for single-exit FedAvg, the latter being poorly trained under the same budget). The multi-exit FedAvg cascade saves less (1.6 to 4.0 times) because its shallow exits are weaker: only 7--48\% of windows leave at the first exit, against 55--98\% for EcoFed. Table~\ref{tab:exits} shows why: with exit weights $\lambda_j\propto j^2$ the first exit receives only 3.3\% of the loss weight of a full-depth client, and FedAvg's first exit reaches 55.9\% on UCI HAR against 85.3\% for EcoFed, whose depth-1 clients train it with the full loss weight. This is a consequence of our choice of exit weights (made to favour the last exit, \ref{app:extra}); with uniform weights the shallow exits of the energy-blind baselines would be better and their last exit worse, so the cascade comparison should be read with that choice in mind. Second, \emph{the extra training cost of a multi-exit model is negligible.} At $\phi=0.25$ on PAMAP2 EcoFed trains with 12~J more than FedAvg (148 against 136~J) but saves 572~$\mu$J per window relative to single-exit FedAvg; by Eq.~\eqref{eq:breakeven} the extra training energy is repaid after about $2.1\times10^4$ windows, i.e.\ about 22 minutes of operation of the 20-device fleet. On UCI HAR and at $\phi=1$ EcoFed trains with \emph{less} energy than FedAvg and also serves with less, so $M^\star\le0$. Third, \emph{serving dominates the lifecycle.} In every row serving matches the training energy within $0.1$ to $1.7$ days of fleet operation (Table~\ref{tab:life}, last column; Fig.~\ref{fig:life}), and over a month of operation inference would use roughly 18 to 270 times the energy of training. Savings in the training phase are therefore worth having mainly because they come with models whose deployed cost is also lower; the largest lever on the total energy is the exit policy.

\begin{figure}[t]
\centering\includegraphics[width=\textwidth]{fig9_lifecycle}
\caption{Lifecycle energy against the number of inferences served by the fleet ($\phi=1$). Dots mark $M_\times$, where inference energy equals training energy.}\label{fig:life}
\end{figure}
